% =============================================================================
%  Migración de GEOUNED a pyOCC — Informe técnico
%
%  Documento autocontenido. Compilar con:
%      pdflatex pyocc_migration_report.tex
%      pdflatex pyocc_migration_report.tex   (segunda pasada: índice + refs)
%
%  Sintetizado a partir del registro de trabajo CLAUDE.md del proyecto.
%  Rama: georeverse-migration. Fecha: 2026-09-06.
% =============================================================================
\documentclass[11pt,a4paper]{article}

\usepackage[utf8]{inputenc}
\usepackage[T1]{fontenc}
\usepackage{lmodern}
\usepackage[spanish]{babel}
\usepackage[a4paper,margin=2.4cm]{geometry}
\usepackage{textcomp}
\usepackage{xcolor}
\usepackage{booktabs}
\usepackage{longtable}
\usepackage{array}
\usepackage{ragged2e}
\usepackage{enumitem}
\usepackage{parskip}
\usepackage{microtype}
\usepackage[hidelinks]{hyperref}
\setlength{\emergencystretch}{3em}

% --- estilo --------------------------------------------------------------------
\definecolor{accent}{HTML}{1F6F8B}
\hypersetup{
  pdftitle={Migracion de GEOUNED a pyOCC - Informe tecnico},
  pdfsubject={Refactor del backend CAD de GEOUNED},
  colorlinks=true, linkcolor=accent, urlcolor=accent, citecolor=accent
}

% Identificadores/rutas literales: \detokenize imprime _, #, &, %, ~ tal cual.
\providecommand{\code}[1]{\texttt{\detokenize{#1}}}
\newcommand{\statusok}{\textbf{\textcolor{accent}{reparado}}}
\newcommand{\statusdone}{\textbf{\textcolor{accent}{resuelto}}}
\newcommand{\statusnone}{\textit{sin reparación}}
\newcommand{\statuslimit}{\textit{limitación aceptada}}

\newcolumntype{L}[1]{>{\RaggedRight\arraybackslash}p{#1}}

\setlist[itemize]{leftmargin=1.4em, itemsep=2pt, topsep=3pt}
\setlist[enumerate]{leftmargin=1.6em, itemsep=4pt, topsep=3pt}

\title{Migración de GEOUNED a pyOCC\\[6pt]
       \large\normalfont Informe técnico del refactor del backend CAD}
\author{Proyecto GEOUNED \textemdash{} rama \texttt{georeverse-migration}}
\date{6 de septiembre de 2026}

% =============================================================================
\begin{document}
\maketitle

\begin{center}
\small
\begin{tabular}{@{}ll@{}}
\toprule
Proyecto            & GEOUNED-org/GEOUNED \\
Rama                & \code{georeverse-migration} \\
Motor por defecto   & \code{ocp} (OCP / pybind11) \\
Motores soportados  & \code{freecad} $\cdot$ \code{occ} $\cdot$ \code{ocp} \\
Alcance             & CadToCsg (GEOUNED) + CsgToCad (GEOReverse) \\
Fuente              & registro de trabajo \code{CLAUDE.md} \\
\bottomrule
\end{tabular}
\end{center}

\vspace{4pt}
\noindent\fbox{\parbox{\dimexpr\textwidth-2\fboxsep-2\fboxrule}{%
\small\centering
\textbf{3\,/\,3} motores CAD con la suite en verde \quad$\vert$\quad
\textbf{0} partículas perdidas (corpus \code{test_models}, 143 ejec.\ d1suned) \quad$\vert$\quad
\textbf{93,6\,\%} tallies dentro de 2$\sigma$ / \textbf{0\,\%} por encima de 3$\sigma$ \quad$\vert$\quad
\textbf{0} imports directos de \code{Part}/\code{FreeCAD} en GEOUNED%
}}

\vspace{6pt}
\tableofcontents
\vspace{6pt}
\hrule

% =============================================================================
\section{Resumen ejecutivo}

GEOUNED convierte geometría CAD en geometría CSG (y a la inversa) para los
códigos de transporte Monte Carlo \textbf{MCNP, OpenMC, PHITS y Serpent}.
Históricamente dependía por completo de la API \code{Part} de FreeCAD y, en
particular, de \code{BOPTools.SplitAPI.slice} para descomponer sólidos. Ese
acoplamiento se ha eliminado.

El trabajo realizado se agrupa en cinco bloques:

\begin{itemize}
  \item \textbf{Aislamiento del kernel.} Todo el código que toca FreeCAD/OCCT
        vive ahora en un único paquete, \code{geo}, tras una interfaz de
        nombres estable. Ningún otro fichero de GEOUNED importa \code{Part} ni
        \code{FreeCAD}.
  \item \textbf{Tres motores geométricos.} \code{geo} tiene tres
        implementaciones intercambiables mediante la variable de entorno
        \code{GEOUNED_CAD_ENGINE}: FreeCAD (\code{freecad}), pythonocc-core /
        SWIG (\code{occ}) y OCP / pybind11 (\code{ocp}, actualmente el valor por
        defecto).
  \item \textbf{Verificación cuantitativa.} Cada corrección se ha validado con
        dos métodos independientes: comprobación punto a punto de
        \code{check_sign} contra geometría CAD construida de forma
        independiente, y un chequeo estocástico de volumen ejecutando el modelo
        MCNP resultante en d1suned (tally $\approx$ 1,0).
  \item \textbf{Depuración geométrica dirigida por datos.} Decenas de bugs
        reales \textemdash{} muchos preexistentes a la migración \textemdash{}
        encontrados y corregidos siguiendo trazas concretas: signos AND/OR en
        superficies compuestas, degeneraciones de tangencia en \code{Gsplit},
        clasificación de superficies, escalado de tolerancias y \emph{slivers}
        residuales.
  \item \textbf{Reparación de defectos CAD.} Una capa nueva en tiempo de carga
        que detecta sólidos corruptos (aristas patológicamente cortas, anillos
        de contorno duplicados, cortes fantasma) y los repara automáticamente
        cuando es rápido y fiable, o los marca sin traducir cuando no lo es.
\end{itemize}

El pipeline inverso, \textbf{GEOReverse} (CsgToCad), que inicialmente quedó
fuera de alcance por su dependencia dura del formato \code{.FCStd}, ha recibido
su propia migración en la rama actual: es engine-switchable, exporta STEP real
vía XCAF bajo los motores pyOCC y colorea cada sólido por material.

Estado a fecha del informe: los tres motores pasan sus suites; el corpus de
referencia \code{Solidos/test_models} se traduce sin partículas perdidas y sin
fallos por encima de 3$\sigma$; \code{occ} y \code{ocp} producen resultados
idénticos sobre todo el corpus. Quedan abiertos varios puntos de GEOReverse y
algunas superficies exóticas sin portar (sección~\ref{sec:pending}).

% =============================================================================
\section{Contexto y motivación}

\subsection{El problema que originó la migración}

El \code{Cut} booleano de FreeCAD puede fallar de forma \textbf{silenciosa} al
partir un sólido en dos cuando la intersección del plano de corte coincide
exactamente con una línea de tangencia preexistente (por ejemplo, un plano que
corta la esquina donde un cilindro ya es tangente a dos caras de un cubo). El
síntoma: \code{Cut} devuelve el sólido original sin cortar, en lugar de lanzar
un error. La causa: los fragmentos resultantes solo se tocan a lo largo de una
línea (contacto de área cero), lo que produce una forma no-manifold que el
algoritmo BOP de OCC no resuelve limpiamente. Pasar por una
exportación/importación STEP «arregla» el problema porque el sanador de STEP
vuelve a partir el compound.

La solución de fondo \textemdash{} construir un grafo de adyacencia de caras
que excluya las aristas compartidas por más de dos caras y reconstruir un
sólido por componente conexa \textemdash{} requiere control topológico fino que
la API \code{Part} de FreeCAD, al ser ella misma un envoltorio sobre OCCT, no
expone cómodamente.

\subsection{Decisión: pythonocc-core, no CadQuery}

\begin{itemize}
  \item Se necesita control topológico de grano fino (grafos de adyacencia
        cara/arista, \code{BRepCheck_Analyzer}, \code{ShapeFix_Shape}, control
        de tolerancia sobre \code{BRepAlgoAPI_*} / \code{BOPAlgo_Splitter}) para
        exactamente los casos degenerados anteriores.
  \item Migrar a pyOCC \emph{elimina} una capa (el envoltorio \code{Part}) en
        lugar de sustituirla por otra distinta (CadQuery).
  \item CadQuery encaja mejor en construcción paramétrica de alto nivel, no en
        la cirugía topológica de bajo nivel que exige la descomposición de
        GEOUNED.
\end{itemize}

% =============================================================================
\section{Arquitectura: el paquete \texorpdfstring{\code{geo}}{geo}}

\subsection{Intento 1 (retirado): la ABC \texorpdfstring{\code{GeometryBackend}}{GeometryBackend}}

El primer diseño fue un patrón Adapter / Ports \& Adapters clásico: una clase
base abstracta \code{GeometryBackend} con dataclasses neutras (\code{GSolid} /
\code{GFace} / \code{GEdge} / \code{GVector} / \dots) que llevaban un objeto
\code{native} más una referencia al backend que las produjo, y una
implementación concreta \code{FreeCADBackend}. Todo GEOUNED llamaba a través de
una única instancia inyectada (\code{_backend.split(gsolid, tool, tol)}, \dots).
Se implementó por completo y se validó contra toda la suite, pero se juzgó
\textbf{demasiado pesado e indirecto} \textemdash{} las dataclasses cargando su
propio backend y la convención \code{_backend.method(x, ...)} \textemdash{} para
lo que aportaba. Se eliminó por completo junto con sus tests.

\subsection{Intento 2 (actual): el paquete \texorpdfstring{\code{geo}}{geo}}

\code{src/geouned/geo/} es el \textbf{único lugar de GEOUNED autorizado a
importar} \code{Part} / \code{FreeCAD} / \code{BOPTools}.

\begin{longtable}{@{}L{0.30\textwidth} L{0.64\textwidth}@{}}
\toprule
\textbf{Fichero} & \textbf{Responsabilidad} \\
\midrule
\endhead
\code{vector_geometry.py} & Matemática pura sin dependencia de FreeCAD:
  \code{GVector}, \code{GMatrix}, \code{GBoundBox} y conversores. \\
\code{surface_geometry.py} & Predicados analíticos sobre descriptores de
  superficie (\code{is_same_*_surface}, \code{is_inside_*},
  \code{find_can_plane}, intersecciones plano-plano y línea-línea). \\
\code{solid_defects.py} & Detección de defectos CAD a nivel de sólido completo
  (\code{find_short_edges}, \code{find_split_ring_faces},
  \code{near_surface_pair}). \\
\code{io_utils.py} & Utilidades de proceso/carga: \code{GLabelNode} (árbol de
  etiquetas STEP), \code{suppress_native_stdout}. \\
\code{_freecad_impl.py} & Implementación FreeCAD: clasificación de
  superficies/curvas, clases de topología
  (\code{GEdge}/\code{GWire}/\code{GFace}/\code{GShell}/\code{GSolid})
  construidas ansiosamente con comportamiento real como métodos, constructores
  \code{Gmake_*} y operaciones
  \code{Gcut}/\code{Gcommon}/\code{Gfuse}/\code{Gsplit}. \\
\code{_occ_impl.py} & La misma superficie de nombres contra pythonocc-core
  (SWIG). \\
\code{_ocp_impl.py} & La misma superficie de nombres contra OCP (pybind11). \\
\code{__init__.py} & Punto único de import. Lee \code{GEOUNED_CAD_ENGINE} una
  sola vez y reexporta los nombres desde la implementación resuelta. \\
\bottomrule
\end{longtable}

\paragraph{Principio de intercambiabilidad.}
La selección de motor ocurre en tiempo de import, no por inyección de
dependencias. Un motor nuevo solo tiene que definir los mismos nombres de clase
y función; \code{geo/__init__.py} elige cuál reexportar y el resto de GEOUNED no
cambia. GEOReverse replica exactamente el mismo patrón con su propio
\code{GEOReverse/Modules/__init__.py}.

\subsection{Diseño de las clases \texorpdfstring{\code{geo}}{geo}}

\begin{itemize}
  \item Las caras se clasifican ansiosamente en una de 5 superficies analíticas
        (plano / cilindro / cono / esfera / toro) vía \code{Gclassify_surface}.
        Las superficies compuestas (RoundCorner, Can, TCone, MultiPlane,
        MultiRoundCorner, ReversedConeCylinder) las ensambla el propio GEOUNED
        con \code{Gcut}/\code{Gcommon}/\code{Gfuse}; \code{geo} nunca las modela.
  \item La clasificación de curvas es tolerante: devuelve \code{None} para una
        arista degenerada o un tipo real no soportado (p.\,ej.\ hipérbola) en
        lugar de lanzar, porque los sólidos se construyen ansiosamente y una
        arista incidental no debe abortar el sólido entero.
  \item \code{Gsplit()} devuelve un \code{SplitResult} que nunca puede volver
        vacío silenciosamente: los casos de tangencia/degeneración se resuelven
        internamente (reintento de tolerancia, reconstrucción no-manifold,
        separación por componentes) y se reportan con
        \code{degenerate_case_handled=True}.
  \item \code{utils/geometry_gu.py::FaceGu} / \code{SolidGu} \textemdash{} un
        segundo envoltorio nativo previo a \code{geo}, usado por el código de
        descomposición \textemdash{} ahora \textbf{heredan} de \code{GFace} /
        \code{GSolid} y solo añaden encima los extras del lado de
        descomposición, en lugar de duplicar la construcción.
\end{itemize}

% =============================================================================
\section{Desacoplamiento de FreeCAD en GEOUNED}

El cierre de los \code{import Part} / \code{import FreeCAD} se hizo fichero por
fichero, trazando cada consumidor real antes de tocar la fuente. \code{core.py}
(concretamente \code{_set_geometry_bounding_box}) fue el último reducto,
precisamente porque su \code{BoundBox} se consume fuera del alcance original de
la migración (\code{void.py}, \code{write_files.py}). Tras su cierre,
\textbf{cero ficheros} del subpaquete GEOUNED importan
\code{Part}/\code{FreeCAD}/\code{BOPTools} directamente.

\subsection{Unificación de representaciones duplicadas}

Para «un plano» existían tres representaciones paralelas: el descriptor
\code{geo.GPlane}, la clase \code{PlaneParams} de \code{basic_functions_part1.py}
y la clase \code{Plane} de \code{build_region/Objects.py}. Se consolidaron:

\begin{itemize}
  \item «superficie analítica + bounding box $\rightarrow$ sólido finito
        recortado» para plano/cilindro/cono se unificó en tres funciones puras
        en \code{build_region/Objects.py} (\code{plane_polygon_from_box},
        \code{cylinder_from_box}, \code{cone_from_box}), en espacio
        \code{GVector}/\code{GBoundBox}, sin duplicación.
  \item Las clases \code{Plane}/\code{Cylinder}/\code{Cone}/\code{Sphere} de
        \code{Objects.py} colapsaron en una única \code{CellSurface} que
        envuelve un descriptor \code{geo}; \code{get_surface()} construye ahora
        \code{GPlane}/\code{GCylinder}/\dots directamente vía
        \code{.from_values(...)}, sin paso de traducción ni desvío nativo.
  \item Los \code{*OnlyParams} de Tier-1 pasaron a almacenar \code{GVector}
        directamente en lugar de \code{FreeCAD.Vector};
        \code{check_sign_primitive} se colapsó a un despacho por diccionario
        contra las funciones \code{is_inside_*} compartidas.
\end{itemize}

\paragraph{Lección metodológica registrada.}
Un estudio estático de «quién lee este campo» es necesario pero no suficiente
cuando el valor pasa por un objeto intermedio antes de llegar a la operación
que le importa el tipo. Varias veces, una encuesta por \code{grep} concluyó que
todos los consumidores toleraban \code{GVector}, y la suite real reveló sitios
nativos que el \code{grep} no vio. Norma: seguir cada
\code{AttributeError}/\code{TypeError} real de la suite completa, no solo el
trazado estático.

\subsection{Dos ámbitos de numeración de superficies}

Cada \code{GeounedSurface} tiene \code{.bVar} (el id firmado con que se
referencia) y, para tipos compuestos, \code{.region} (la definición booleana
AND/OR) y \code{.components} (la relación numeración\,$\leftrightarrow$\,superficie).
Se asignan en dos fases con numeraciones genuinamente distintas: en
\textbf{descomposición} el \code{.bVar} es un id local desechable cuyo único fin
es mantener la expresión AND/OR consistente para una construcción CAD; en
\textbf{conversión} se sobreescribe con el id global canónico y deduplicado.
Ambas fases \emph{no} pueden leer la \code{.region} de la otra, pero sí
comparten la \textbf{regla} AND/OR, extraída a funciones puras
(\code{can_region}, \code{tcone_region}, \code{round_corner_region}, \dots).

% =============================================================================
\section{Reorganización de los writers}

Los cuatro escritores de formato de salida tenían implementaciones casi
idénticas de varios métodos. Se reorganizaron por la frontera real de
convenciones de tarjeta:

\begin{itemize}
  \item \code{write/mcnp_like/} \textemdash{} MCNP, Serpent y PHITS (los tres
        formatos de texto de linaje MCNP, con convenciones comunes de comentario
        de línea/inline y banners).
  \item \code{write/openmc/} \textemdash{} estructuralmente distinto (salida
        XML/Python, sin concepto de tarjeta de comentario).
  \item \code{write/functions.py} \textemdash{} hogar compartido de
        \code{get_cell_surf_summary}, \code{simplify_planes},
        \code{sorted_surfaces}, usadas por los cuatro formatos.
\end{itemize}

El mixin \code{CommonInputWriter} (heredado por
\code{McnpInput}/\code{SerpentInput}/\code{PhitsInput}) recoge lo duplicado solo
entre los tres formatos de linaje MCNP, parametrizado por atributos de clase
(\code{_surface_formatter}, \code{inline_comment_char}, \dots).

\paragraph{Corrección real, confirmada con el usuario.}
La unificación de \code{sorted_surfaces} corrige una inconsistencia: solo MCNP
aplicaba \code{Surfaces.IndexOffset} al reetiquetar cada \code{bVar};
Serpent/OpenMC/PHITS se lo saltaban. Ahora los cuatro formatos usan la versión
de MCNP. El efecto es actualmente latente (\code{IndexOffset} es 0 salvo que
\code{settings.startSurf} se cablee, cosa que también se hizo en esta migración)
pero deja de ser una divergencia entre formatos.

% =============================================================================
\section{Campañas de verificación}
\label{sec:verif}

\subsection{Verificación de \texorpdfstring{\texttt{check\_sign}}{check\_sign} punto a punto}

Para cada superficie compuesta encontrada en un modelo, se construye un sólido
CAD de referencia \emph{independiente} por la misma vía que usa \code{Gsplit}
para construir esa composición, se muestrean puntos aleatorios estrictamente
dentro de la caja de construcción y se compara \code{check_sign(punto,
superficie)} contra \code{gsolid.is_inside(punto)}.

\begin{longtable}{@{}L{0.22\textwidth} L{0.15\textwidth} L{0.57\textwidth}@{}}
\toprule
\textbf{Tipo compuesto} & \textbf{Resultado} & \textbf{Notas} \\
\midrule
\endhead
RoundCorner & 300/300 & Múltiples instancias en decenas de ficheros. \\
Can (Fwd / Rev) & 300/300 & Un artefacto de muestreo (padding de bounding box)
  se identificó y corrigió: muestrear dentro de la caja de construcción, no de
  la del sólido resultante. \\
TCone (Fwd / Rev) & 300/300 & Reveló un bug real en \code{TCone_region}
  (ignoraba \code{p1_configuration}/\code{p2_configuration}), corregido. \\
MultiRoundCorner & 300/300 & 17 instancias en 13 ficheros. Primera validación
  end-to-end de este tipo. \\
MultiPlane & 300/300 & Reveló que \code{multiplane()} nunca podía detectar un
  MultiPlane (comparación instancia-vs-clase), corregido. \\
ReversedConeCylinder & parcial & Estructuralmente inalcanzable desde el escaneo
  de descomposición; verificado desde el lado de conversión con una discrepancia
  real en \code{hylife-v06} solid113 (271/300), ligada a un desajuste conocido
  de GEOReverse. \\
\bottomrule
\end{longtable}

\subsection{Chequeo estocástico de volumen MCNP}

Pregunta distinta y más end-to-end: ¿un modelo MCNP \emph{completamente
traducido}, ejecutado de verdad, reproduce los \textbf{mismos volúmenes} que los
sólidos CAD originales? Se exporta con \code{volSDEF=True} (fuente fotónica
isótropa desde una esfera envolvente + tally F4 de longitud de traza en cada
celda, normalizado por el volumen CAD), se corre en \code{MODE P} + \code{VOID}
en d1suned y el resultado esperado del tally es exactamente \textbf{1,0}. Una
desviación significativa ($>$~3$\sigma$) señala que los signos de superficie
escritos no reconstruyen el mismo sólido.

\begin{center}
\small
\begin{tabular}{@{}L{0.68\textwidth} r@{}}
\toprule
\textbf{Corpus \code{Solidos/test_models} (sin \code{Big_*})} & \textbf{Valor} \\
\midrule
Ficheros que traducen                                  & 133 / 138 \\
Tallies dentro de 2$\sigma$ de 1,0                     & 93,6\,\% \\
Tallies marginales (2\textendash 3$\sigma$, ruido esperado) & 6,4\,\% \\
Fallos reales ($>$ 3$\sigma$)                          & 0\,\% \\
Ficheros con partículas perdidas                       & 0 / 143 ejec. \\
Divergencia \code{occ} vs \code{ocp} sobre el corpus   & byte a byte idéntica \\
\bottomrule
\end{tabular}
\end{center}

Los 5 ficheros que no traducen son defectos CAD genuinos que la capa de
detección (sección~\ref{sec:defects}) marca correctamente, más un caso de crash
nativo aceptado como limitación (\code{ConeSphere.stp} bajo
\code{occ}/\code{ocp}).

% =============================================================================
\section{Correcciones de geometría}

Decenas de bugs reales, la mayoría preexistentes a la migración (confirmado
contra el último commit pre-migración), encontrados siguiendo trazas concretas
\textemdash{} nunca por conjetura \textemdash{} y verificados antes y después.
Agrupados por naturaleza:

\begin{longtable}{@{}L{0.16\textwidth} L{0.56\textwidth} L{0.22\textwidth}@{}}
\toprule
\textbf{Categoría} & \textbf{Ejemplos representativos} & \textbf{Impacto} \\
\midrule
\endhead

\textbf{Signos AND/OR en superficies compuestas} &
\code{can_region} rama cono + plano-ápice Reversed usaba AND donde el signo
físico exige OR (dos instancias); \code{TCone_region} sustituía la configuración
real por un valor único según orientación; \code{round_corner_region} aplicaba
\code{fwd_cyl} dos veces en la rama \code{p1 == p2}; \code{cone_apex_plane}
doble-compensaba el caso Reversed; \code{add_reversedCC} rediseñado
(\code{plane_region AND (s1 AND ... AND sn)}) y su \code{AddPlanes} pasó de OR
incondicional a AND; \code{AdjacentMultiplanePlanes} pasó a lista-de-listas (AND
dentro de grupo, OR entre grupos). &
Volumen CSG erróneo con expresión booleana válida; en varios casos «fugas» que
engullían medio modelo. Verificado con d1suned tras cada corrección. \\
\addlinespace

\textbf{Degeneración / tangencia en \code{Gsplit}} &
Fallback de corte coaxial cono-cono (\code{_try_coaxial_cone_split}):
reconstrucción analítica del círculo exacto de intersección, pre-split de esa
arista y reintento con escalado de tolerancia difusa; el \code{splitTolerance}
por defecto implícito pasó de 0,0 a $10^{-4}$ bajo \code{occ}/\code{ocp} (OCCT
7.9.3 es más sensible a la near-tangencia); \code{_repair_non_manifold_solid}
ahora valida su propia reconstrucción (validez + conservación de volumen) antes
de fiarse de ella. &
Descomposición que devolvía silenciosamente 1 pieza en vez de 2 o 3; fragmentos
con volumen inflado. \\
\addlinespace

\textbf{Clasificación de superficies} &
\code{is_same_plane_surface} comparaba el offset de dos planos con ejes
antiparalelos como \code{d1 == d2} en vez de \code{d1 == -d2};
\code{convex_planes} tenía un test de signo muerto (\code{cross.dot(ref)}
siempre 0) y no normalizaba \code{atan2} a $[0, 2\pi)$; \code{multiplane()}
comparaba instancia contra clase (\code{Gclassify_curve(e) is GLine}), haciendo
MultiPlane indetectable; \code{is_closed_cylinder_cone} aceptaba una cara que no
era un cilindro cerrado de 360$^\circ$ $\rightarrow$ invariante topológica real
por winding (\code{_is_closed_by_winding}); \code{cyl_plane_region_conf} usaba la
cara semilla para ambos extremos de un cilindro partido. &
Clasificación de caras a tipo compuesto equivocado; tallies de 0,0 y
desviaciones de cientos de $\sigma$, todas cerradas. \\
\addlinespace

\textbf{Tolerancias} &
\code{min_area} nunca estaba cableado a \code{order_plane_face} ni a los
generadores de candidatos; nuevo \code{min_face_width} + \code{CharacteristicWidth}
(ancho característico de una cara vía radios de giro) para detectar
\emph{slivers} «estrechos pero de área suficiente», portado a los tres motores;
\code{Tolerances.scaled(volume)} escala \code{min_area}/\code{min_face_width}
hacia abajo para fragmentos descompuestos diminutos; \code{eligible_plane}
ignoraba las tolerancias configuradas por el usuario. &
Planos sliver espurios entrando en la definición de celda $\rightarrow$
partículas perdidas. \\
\addlinespace

\textbf{Slivers residuales de corte booleano} &
\code{other_face_edge} con modo \code{skip_slivers} para atravesar caras sliver
y encontrar la cara real; filtro de slivers en \code{merge_same_surface_faces}
antes del walk $O(n^2)$ de adyacencia; \code{sort_range} con fallback «sin
progreso» ante fragmentos unidos solo módulo $2\pi$ (era recursión infinita). &
Detección de Can/RoundCorner fallando o colgándose sobre geometría con
artefactos de corte. \\
\addlinespace

\textbf{Cross-engine} &
El flag de \code{ShapeUpgrade_UnifySameDomain} que provoca un crash nativo es
\emph{opuesto} entre OCP y pythonocc-core sobre la misma geometría OCCT 7.9.3;
\code{get_join_cone_cyl} usaba \code{.Index} donde un shell fusionado tiene
\code{.Indexes} (solo \code{freecad}); el LAPACK de MKL entraba en conflicto con
\code{occt}/\code{tbb} en el mismo proceso $\rightarrow$ cambio a OpenBLAS en
\code{pyoccenv} y \code{ocpenv}. &
47 tests crasheando en \code{freecad}/\code{occ} por un \code{AttributeError};
access violations nativos. \\
\addlinespace

\textbf{No geométricos, encontrados de camino} &
\code{generic_split} hacía doble-wrap de un \code{GSolid} ya envuelto
$\rightarrow$ la descomposición era silenciosamente un no-op;
\code{check_intersection} con early-exit por \code{Gdistance} clasificaba
sólidos totalmente anidados como disjuntos $\rightarrow$ la celda de vacío del
enclosure no excluía las esferas dentro (\code{w_encl.stp}, tally 0,0);
\code{expand_regions_to_integer} sin guard de \code{bool}, guard que los cinco
escritores de celda necesitaban y solo MCNP tenía. &
Bugs de correctitud que degradaban la descomposición o la salida de forma
invisible. \\
\bottomrule
\end{longtable}

% =============================================================================
\section{Migración pyOCC por fases}

La migración a pyOCC se abordó en fases, cada una con un criterio de
continuación explícito antes de escribir código a gran escala.

\begin{enumerate}
  \item \textbf{Fase 1 \textemdash{} Scaffolding y prueba go/no-go.} Selección
        de motor vía \code{GEOUNED_CAD_ENGINE}; \code{geo/__init__.py} importa de
        \code{_freecad_impl} o del nuevo \code{_occ_impl}. Solo se implementó lo
        mínimo para probar \code{Gsplit} contra el caso \code{rev_pipe.stp}, ya
        caracterizado como roto en FreeCAD. Resultado: \code{BOPAlgo_Splitter}
        (OCCT 7.9.3) devuelve \textbf{3 sólidos válidos} en el primer intento
        sin reparación, donde FreeCAD devolvía 1 sólido inválido. Señal fuerte
        para continuar.
  \item \textbf{Fase 2 \textemdash{} Port completo.} Los 5 descriptores de
        superficie y 4 de curva, \code{Gclassify_surface}/\code{Gclassify_curve},
        \code{GEdge}/\code{GWire}/\code{GFace}/\code{GShell}/\code{GSolid}
        completos (incluyendo \code{find_interior_point},
        \code{faces_sharing_edge}, \code{fix}/\code{refine}), todos los
        \code{Gmake_*} y las operaciones booleanas. Test nuevo
        \code{tests/geo/test_occ_impl.py}: 39/39 contra geometría real con
        valores esperados conocidos.
  \item \textbf{Fase 3 \textemdash{} Pipeline de integración.}
        \code{tests/test_cadtocsg.py} pasa 50/50 bajo
        \code{GEOUNED_CAD_ENGINE=occ}, idéntico al baseline FreeCAD, tras cerrar
        6 huecos reales (entre ellos \code{Gload_step_labels} vía XCAF, un bug de
        periodicidad en \code{get_join_cone_cyl} alcanzable también bajo FreeCAD,
        y el fallback de \code{GeomAPI_ProjectPointOnSurf}).
  \item \textbf{Fase 4 \textemdash{} Corpus \code{Solidos/}.} Escaneo diferencial
        de 98/99 ficheros contra un baseline FreeCAD fresco; 7 huecos más
        cerrados. \code{rev_pipe.stp} confirmado: FreeCAD falla la ruta completa
        decompose+convert (\code{FuseEdges}), \code{occ} tiene éxito con
        \code{RoundC=2, MultiRoundC=1}.
\end{enumerate}

\subsection{Tercer motor: OCP (pybind11)}

Añadido junto a FreeCAD y pythonocc-core, no como reemplazo. OCP es un binding
pybind11 de OCCT más directo y activamente mantenido que el SWIG de
pythonocc-core. Validación go/no-go contra \code{rev_pipe.stp} (2 piezas,
volumen conservado), port completo estructuralmente copiado de
\code{_occ_impl.py} y retargeteado (con diferencias de API confirmadas en vivo:
sufijo \code{_s} en métodos estáticos, sin \code{.DownCast()}, orden de
argumentos invertido en \code{ShapeUpgrade_UnifySameDomain}, \dots).
\code{tests/geo/test_ocp_impl.py} 39/39 al primer intento;
\code{tests/test_cadtocsg.py} 50/50; \code{tests/test_csgtocad.py} 2/2. Ambos
motores pyOCC se desarrollan en paralelo: una corrección en uno se porta al
otro.

\paragraph{El «gap» de rendimiento, contado con honestidad.}
En el caso motivador (\code{hylife-v06.stp} solid 17), FreeCAD tardaba 2,3\,s y
pythonocc-core 19,2\,s. Se atribuyó primero al coste por-llamada del binding
SWIG (real: pythonocc-core expone la API cruda de OCCT sin la capa de
conveniencia en C++ que FreeCAD sí tiene). OCP en ese mismo caso: 18,4\,s
\textemdash{} solo $\sim$4\,\% más rápido, no el 2\textendash 6$\times$ de los
micro-benchmarks de conversión de puntos. La causa real resultó ser
\textbf{otra}: \code{splitTolerance} con valor implícito 0,0 hacía que OCCT
7.9.3 computara y luego descartara un split degenerado de 15 fragmentos. Con el
default corregido a $10^{-4}$ bajo \code{occ}/\code{ocp}, el mismo caso baja a
\textbf{1,2\,s} \textemdash{} más rápido que FreeCAD.

% =============================================================================
\section{Migración de GEOReverse}

GEOReverse (CsgToCad) reconstruye CAD a partir de CSG y depende históricamente
del formato de documento \code{.FCStd} de FreeCAD, no solo de su kernel
geométrico. Inicialmente quedó permanentemente fuera de alcance; en la rama
actual (\code{georeverse-migration}) ha recibido el mismo tratamiento que el
pipeline directo.

\begin{itemize}
  \item \textbf{Desacoplamiento.} El núcleo de GEOReverse (\code{Objects.py},
        \code{buildSolidCell.py}, \code{splitFunction.py}, \code{buildCAD.py},
        \code{MCNPinput.py}) importa \code{geo} directamente, igual que el
        pipeline directo. Un módulo \code{matrix_utils.py} recoge el puente
        numpy\,$\leftrightarrow$\,matriz nativa.
        \code{BOPTools.SplitAPI.slice} $\rightarrow$ \code{Gsplit};
        \code{FuseSolid} (triplicado byte a byte) $\rightarrow$ un único
        \code{fuse_solids} con la comprobación de invariancia de volumen de
        \code{GSolid.refine()}.
  \item \textbf{\code{export_cad} independiente del motor.}
        \code{export_cad(output_filename="", format="stp")} acepta un \code{str}
        o una lista de formatos, valida contra los formatos soportados por el
        motor y despacha por \code{CAD_ENGINE}. Cambio de comportamiento real:
        el \code{.FCStd} ya no se escribe incondicionalmente, solo cuando se
        pide \code{"fcstd"} explícitamente.
  \item \textbf{Simetría FreeCAD/pyOCC completa.}
        \code{GEOReverse/Modules/__init__.py} despacha entre
        \code{_freecad_impl.py} (export vía documento FreeCAD, \code{makeTree}) y
        \code{_ocp_impl.py}/\code{_occ_impl.py} (export real vía XCAF:
        \code{STEPCAFControl_Writer}, jerarquía Universe$\rightarrow$Material$\rightarrow$Cell
        con nombres y anidamiento). Las 6 superficies cuádricas exóticas se
        reorganizaron en un paquete \code{geo_quadrics} con el mismo patrón de
        despacho.
  \item \textbf{Color por material (motor \code{occ}/\code{ocp}).} Cada sólido
        del STEP exportado se colorea por su valor \code{MAT} vía XCAF: mismo
        color para el mismo material, paleta \code{tab10} hasta 10 materiales y
        rotación por ángulo áureo para más, sin colisiones. El color en el motor
        FreeCAD queda fuera de alcance (su API de color vive en
        \code{Gui.ViewProvider}, sin ruta headless).
  \item \textbf{IGES descartado.} Se verificó en vivo que
        \code{IGESControl_Writer} no preserva sólidos a través de IGES en esta
        build de OCCT (un sólido vuelve como caras planas recortadas sin
        topología de shell). Se abandonó importación y exportación IGES.
\end{itemize}

Verificación: \code{tests/test_csgtocad.py} es engine-aware y pasa 2/2 bajo los
tres motores contra aserciones de volumen reales por sólido ($\pm 10^{-6}$).
Bajo \code{occ}/\code{ocp} quedan 2 fallos preexistentes (GEOReverse produce
2/1 sólidos donde se esperan 4/5 en el STEP de \code{cylinder_box}), confirmados
no introducidos por esta migración y no ligados al trabajo de \code{Gsplit}.

% =============================================================================
\section{Detección y reparación de defectos CAD}
\label{sec:defects}

Una capa nueva en tiempo de carga (\code{Gload_and_process_step} $\rightarrow$
\code{Gcheck_and_repair}) que, para cada sólido: aplica un \code{fix(1e-6)},
comprueba defectos conocidos y, si los hay, intenta repararlos con el método más
apropiado; solo si la reparación falla se consulta el modo elegido por el
usuario. El parámetro \code{corrupted_solids} tiene dos modos, \code{"stop"}
(por defecto) y \code{"remove"}; el antiguo parámetro \code{invalid_solids} se
fusionó aquí.

\begin{longtable}{@{}L{0.19\textwidth} L{0.29\textwidth} L{0.28\textwidth} L{0.15\textwidth}@{}}
\toprule
\textbf{Defecto} & \textbf{Detección} & \textbf{Reparación} & \textbf{Estado} \\
\midrule
\endhead

\textbf{Anillo de contorno partido} (micro-recorte duplicado / escalón
colapsado, variante esférica) &
\code{find_split_ring_faces}: cara analítica curva a escala sliver que además
toca una arista patológicamente corta. \code{count_split_ring_pairs} cuenta
pares de aristas circulares casi coincidentes y concéntricas. &
\code{Gcollapse_split_rings} (ocp/occ): quita las caras «riser», re-cose el
shell, rehace el sólido. Aceptación por 3 comprobaciones: validez +
$|\Delta V|/V < 3\times10^{-4}$ + los pares de anillos deben decrecer. &
\statusok\newline \code{barrel bottom.stp} $\rightarrow$ tally 0,9997 \\
\addlinespace

\textbf{Escalón cilíndrico colapsado} (familia beltline) &
Mismo detector de anillo partido dispara; \code{near_surface_pair} identifica
dos cilindros de radio casi igual bridged por una banda sliver. &
\code{Gsliver_heal} (ocp/occ), versión 0: quita las caras sliver y la menor del
par cercano, re-recorta la cuádrica liberada a su radio real, re-cose
\emph{al final}. Gate: validez + $|\Delta V| < 5\times10^{-4}$. &
\statusok\newline \code{LR.stp} $\rightarrow$ tally 0,9989 (era 0,0 / 24
partículas perdidas) \\
\addlinespace

\textbf{Corte fantasma} (dos regiones volumétricas que solo se tocan por
aristas \textemdash{} la motivación fundacional de la migración) &
\code{_separate_edge_joined_components}: aristas no-manifold (compartidas por
$\neq$ 2 caras) $\rightarrow$ grafo de adyacencia de caras solo por aristas
manifold $\rightarrow$ si hay $\geq$ 2 componentes conexas y las aristas
no-manifold eran la conexión portante. &
Cose cada componente por separado (sin caras donantes: un corte fantasma real
ya lleva el contorno cerrado de cada región). Gate: $\geq$ 2 piezas, todas
válidas, suma de volúmenes conservada a $10^{-6}$. &
\statusdone\newline ocp/occ $\cdot$ \code{Piece_1}+\code{Tool_1} $\rightarrow$ 2
sólidos válidos (era 1 sólido inválido fusionado). FreeCAD no puede hacerlo. \\
\addlinespace

\textbf{Wire de cara malformado} (\code{UnorientableShape}) &
\code{check_solid_defects} reporta topología inválida real;
\code{BRepCheck_Analyzer} marca la cara, ninguna arista sliver. &
Ninguna. Se probaron \code{ShapeFix_Shape}, \code{ShapeFix_Face}, partición del
wire y reconstrucción \textemdash{} todas fallan o desvían el volumen $>$ 1,7\,\%. &
\statusnone\newline \code{L4_support.stp} movido a
\code{Solidos/Detected_corrupted/}, excluido de conversión \\
\bottomrule
\end{longtable}

Un cambio relacionado: \code{check_solid_defects} ya no rechaza un sólido por
una arista sliver sobre una cara \emph{bien formada} (\code{CharacteristicWidth}
normal) \textemdash{} el pipeline de descomposición ya tolera aristas sliver en
caras sanas, y ningún reparador quita tal arista sin una cara sliver donde
anclarse. Las recetas de defectos se mantienen en un catálogo vivo
(\code{reference_cad_defect_recipes.md}): síntoma observable, firma de detección
general, causa raíz, reparación si es rápida y fiable, y un fixture STEP mínimo
con verificación d1suned.

% =============================================================================
\section{Estado actual y trabajo pendiente}
\label{sec:pending}

\subsection{Estado consolidado}

\begin{itemize}
  \item Los tres motores pasan sus suites: \code{freecad} 158/158, \code{occ}
        128/128, \code{ocp} 128/128 (+2 fallos preexistentes de GEOReverse bajo
        los motores pyOCC).
  \item Corpus \code{Solidos/test_models}: 93,6\,\% $<$ 2$\sigma$, 0\,\% $>$
        3$\sigma$, 0 partículas perdidas en 143 ejecuciones d1suned. \code{occ}
        y \code{ocp} byte a byte equivalentes sobre todo el corpus.
  \item \code{ocp} es el motor por defecto cuando \code{GEOUNED_CAD_ENGINE} no
        está definida.
\end{itemize}

\subsection{Abierto}

\begin{longtable}{@{}L{0.62\textwidth} L{0.32\textwidth}@{}}
\toprule
\textbf{Ítem} & \textbf{Naturaleza} \\
\midrule
\endhead
Discrepancia de volumen en el round-trip de GEOReverse sobre
\code{hylife-v06.stp} celda 1 (CAD reconstruido 1,401$\times$ vs tally d1suned
1,119$\times$). & GEOReverse, no trazado hasta el fondo. \\
\code{test_csgtocad.py} 2 fallos bajo \code{occ}/\code{ocp} (2/1 sólidos donde
se esperan 4/5). & GEOReverse, preexistente, acotado a
\code{Objects.py::CellObj.buildShape}. \\
\code{AdjacentMultiplanePlanes} necesita extenderse de RevCC a
MultiRoundCorner. & Mejora de cobertura, pendiente desde 2026-08-21. \\
Las 6 cuádricas exóticas (\code{Gmake_elliptic_cone}, \code{Gmake_hyperboloid},
\dots) son stubs en los backends pyOCC de GEOReverse. & Port pendiente; sin
fixture que las ejercite todavía. \\
\code{Big_complex_cell/modelcell_cut1.stp} y \code{modelCell_670000.stp} pierden
partículas. & No trazado; excluidos de los escaneos de corpus por convención. \\
\code{hylife-v06.stp} round-trip completo de 372 sólidos: descomposición lenta
en \code{meta_list[45]} (fragmentos de cara duplicados, coste $O(n^2)$ en
\code{same_faces}). & Causa raíz identificada, no optimizado, despriorizado por
el usuario. \\
\code{ConeSphere.stp}: crash nativo (access violation) bajo \code{occ}/\code{ocp}
en \code{ShapeUpgrade_UnifySameDomain}; traduce limpio bajo \code{freecad}. &
\statuslimit{} Sin knob de OCCT que lo evite; documentado en los docstrings. \\
Reorganización y dedup del árbol de fixtures \code{Solidos/}. & Parcial
(\code{test_models/} + algunos \code{duplicates_removed/}). \\
\bottomrule
\end{longtable}

Prioridad explícita del usuario: terminar de limpiar los bugs conocidos del
pipeline directo (GEOUNED) antes de retomar los ítems de GEOReverse.

% =============================================================================
\section{Anexo: entorno y herramientas}

\begin{center}
\small
\begin{tabular}{@{}L{0.20\textwidth} L{0.24\textwidth} L{0.18\textwidth} L{0.30\textwidth}@{}}
\toprule
\textbf{Entorno conda} & \textbf{Binding} & \textbf{OCCT $\cdot$ Python} & \textbf{Uso} \\
\midrule
(python por defecto) & FreeCAD 1.1.1 & 7.8.1 $\cdot$ 3.11 & \code{GEOUNED_CAD_ENGINE=freecad} \\
\code{pyoccenv} & pythonocc-core (SWIG) & 7.9.3 $\cdot$ 3.12 & \code{GEOUNED_CAD_ENGINE=occ} \\
\code{ocpenv} & OCP (pybind11) & 7.9.3 $\cdot$ 3.12 & \code{GEOUNED_CAD_ENGINE=ocp} (por defecto) \\
\bottomrule
\end{tabular}
\end{center}

\begin{itemize}
  \item \textbf{FreeCAD y pyOCC no coexisten en un mismo proceso}
        (\code{Module use of python311.dll conflicts}): FreeCAD trae Python
        3.11, los entornos pyOCC usan 3.12. Cualquier prueba del lado pyOCC usa
        el \code{python.exe} del entorno correspondiente y nunca mezcla en la
        misma invocación.
  \item Los entornos pyOCC necesitaron cambiar el proveedor de BLAS/LAPACK de
        MKL a OpenBLAS
        (\code{conda install -c conda-forge --override-channels "libblas=*=*openblas" ...}):
        el LAPACK de MKL entra en conflicto con \code{occt}/\code{tbb} en el
        mismo proceso y provoca un access violation en \code{numpy.linalg}.
  \item El chequeo de volumen MCNP se ejecuta con \textbf{d1suned} (variante
        MCNP interna) a través de WSL. El pipeline completo, scripts y
        resultados están archivados bajo el WSL (Ubuntu-22.04) en
        \code{~/work/taller/SolidTestMCNP/scripts/} (conversión y d1suned en
        paralelo por motor: \code{run_all_conversions_*_parallel.py},
        \code{run_test_models_*_d1suned.sh}, \code{analyze_test_models_*.py}).
  \item Convención del árbol de trabajo: \code{Solidos/} solo STP;
        \code{Solidos/test_models/} el set curado de regresión;
        \code{../calculos_CLAUDE/} toda ejecución de cálculo;
        \code{Solidos/BadCAD_decomposition/}, \code{Solidos/Detected_corrupted/},
        \code{Solidos/no_convierte/}, \code{Solidos/convierte_bad_volume/}
        carpetas de triaje.
\end{itemize}

\vfill
\begin{center}\footnotesize\itshape
Informe sintetizado a partir del registro de trabajo \code{CLAUDE.md} del
proyecto GEOUNED. Rama \code{georeverse-migration}. Fecha: 2026-09-06.\\
Preferencia de código del proyecto: la comunicación es en español; todo el
código, comentarios, docstrings y nombres de identificadores en inglés.
\end{center}

\end{document}
